\subsection{FIPA ACL: the agent communication language}
\label{sec:fipa-acl}
\slidesrc{Communication: 9--18}

\textbf{FIPA} (Foundation for Intelligent Physical Agents, Section~\ref{sec:research-directions})
publishes the standard specifications for agent platforms (slide 9 lists them by subject:
abstract architecture, agent management, message transport, agent communication,
applications). The \textbf{agent communication} specifications (slide 10) have three parts:

\begin{center}
\begin{forest}
for tree={draw, thick, rounded corners, font=\small, align=center, l sep=6mm, s sep=6mm, edge={thick}}
[{Agent\\Communication}, fill=yellow!40
  [{Interaction\\Protocols}, fill=nbblue!8]
  [{Communicative\\Acts}, fill=nbblue!8]
  [{Content\\Languages}, fill=nbblue!8]]
\end{forest}
\captionof{figure}{FIPA communication specifications (slide 10): \emph{what} kind of act a
message performs (communicative acts), \emph{how} its content is written (content languages),
and \emph{which sequences} of messages are allowed (interaction protocols).}
\end{center}

\subsubsection{Message structure}
\slidesrc{Communication: 11}

\begin{nbfact}{FIPA ACL message parameters (slide 11)}
\begin{center}
\small
\begin{tabularx}{\linewidth}{@{}>{\raggedright\arraybackslash}p{3.3cm} l X@{}}
\toprule
\textbf{Category} & \textbf{Parameter} & \textbf{Meaning (added)} \\
\midrule
type of communicative act & \code{performative} & which act: \code{inform}, \code{request}, \ldots \\
\midrule
participants in communication & \code{sender} & who sends \\
 & \code{receiver} & who receives (one or more agents) \\
 & \code{reply-to} & where replies should go (if not to the sender) \\
\midrule
content of message & \code{content} & what the act is about \\
\midrule
description of content & \code{language} & language of the content (e.g.\ \code{fipa-sl}) \\
 & \code{encoding} & encoding of the content \\
 & \code{ontology} & vocabulary giving meaning to the symbols \\
\midrule
control of conversation & \code{protocol} & interaction protocol being followed \\
 & \code{conversation-id} & identifies the conversation \\
 & \code{reply-with} & label the answer must refer to \\
 & \code{in-reply-to} & label of the message being answered \\
 & \code{reply-by} & deadline for the answer \\
\bottomrule
\end{tabularx}
\end{center}
Only \code{performative} is mandatory; most messages also have sender, receiver and content.
\end{nbfact}

\begin{nbexample}{a request from a buyer to a seller}
\notinslides
\begin{minted}[linenos=false]{text}
(request
   :sender          (agent-identifier :name buyer@shop-platform)
   :receiver        (set (agent-identifier :name seller@shop-platform))
   :content         "((action (agent-identifier :name seller)
                              (deliver book42 (address Genova))))"
   :language        fipa-sl
   :ontology        book-trading
   :protocol        fipa-request
   :conversation-id order-7781
   :reply-with      req-1
   :reply-by        20261110T120000000Z)
\end{minted}
The seller's answer \code{agree} will carry \code{:in-reply-to req-1} and the same
\code{:conversation-id}, so the buyer can match it with its question even if it is running
many conversations at once.
\end{nbexample}

\subsubsection{Communicative acts}
\slidesrc{Communication: 12--13}

\begin{center}
\small
\renewcommand{\arraystretch}{1.12}
\begin{tabularx}{\linewidth}{@{}l X@{}}
\toprule
\textbf{Communicative act} & \textbf{Meaning (from the FIPA Communicative Act Library, added)} \\
\midrule
\code{accept-proposal} & accept a previously submitted proposal to perform an action \\
\code{agree} & agree to perform some action, possibly in the future \\
\code{cancel} & tell another agent that the sender no longer intends it to perform an action requested earlier \\
\code{cfp} (call for proposal) & call for proposals to perform a given action \\
\code{confirm} & tell the receiver that a proposition is true, when it is uncertain about it \\
\code{disconfirm} & tell the receiver that a proposition is false, when it believes or suspects it is true \\
\code{failure} & tell that an action was attempted but failed \\
\code{inform} & tell that a proposition is true \\
\code{inform-if} & tell \emph{whether} a proposition is true \\
\code{inform-ref} & tell the object that corresponds to a description (e.g.\ ``the price is 12'') \\
\code{not-understood} & the sender did not understand a message \\
\code{propagate} & the receiver must forward the embedded message to the agents matching a description, and treat it as sent to itself \\
\code{propose} & submit a proposal to perform an action, under certain conditions \\
\code{proxy} & the receiver must forward the embedded message to the agents matching a description \\
\code{query-if} & ask whether a proposition is true \\
\code{query-ref} & ask for the object referred to by an expression \\
\code{refuse} & refuse to perform an action, giving the reason \\
\code{reject-proposal} & reject a proposal during a negotiation \\
\code{request} & ask the receiver to perform an action \\
\code{request-when} & ask to perform an action when a proposition becomes true \\
\code{request-whenever} & ask to perform an action every time a proposition becomes true \\
\code{subscribe} & ask to be notified of the value of a reference, and again whenever it changes \\
\bottomrule
\end{tabularx}
\captionof{table}{The 22 FIPA communicative acts listed on slides 12--13.}
\end{center}

\begin{nbintuition}[Families of acts]
\emph{Information passing}: inform, inform-if, inform-ref, confirm, disconfirm.
\emph{Requesting information}: query-if, query-ref, subscribe. \emph{Performing actions}:
request, request-when, request-whenever, agree, refuse, cancel, failure.
\emph{Negotiation}: cfp, propose, accept-proposal, reject-proposal. \emph{Error handling}:
not-understood, failure. \emph{Forwarding}: propagate, proxy.
\end{nbintuition}

\subsubsection{Anatomy of a communicative act: accept-proposal}
\slidesrc{Communication: 14--18}

Each act is specified by a summary, its content, a description, a \textbf{formal model} and an
example. The slides show \code{accept-proposal}.

\begin{itemize}
  \item \textbf{Summary} (slide 14): the action of accepting a previously submitted proposal
        to perform an action.
  \item \textbf{Content} (slide 15): a tuple of an \emph{action expression} (the action to be
        done) and a \emph{proposition} (the conditions of the agreement).
  \item \textbf{Description} (slide 16): a general-purpose acceptance of a proposal that was
        previously submitted (typically through a \code{propose} act). The agent sending the
        acceptance informs the receiver that it intends that (at some point in the future)
        the receiver will perform the action, once the given precondition is, or becomes,
        true. The proposition indicates the preconditions the agent is attaching to the
        acceptance: e.g.\ after an offer to ``hold a meeting anytime on Tuesday'', the
        acceptance may add the condition that the meeting is at 11.00. (Future extension: an
        agent may intend that an action is done without necessarily intending the
        precondition, e.g.\ ``I bid price $p$, but I am prepared to accept $p'$''.)
\end{itemize}

\begin{nbdef}{Formal model of accept-proposal (slide 17)}
\[
\langle i, \mathit{accept\text{-}proposal}(j, \langle j, act\rangle, \phi)\rangle \;\equiv\;
\langle i, \mathit{inform}(j, I_i\,\mathit{Done}(\langle j, act\rangle, \phi))\rangle
\]
\begin{align*}
\text{FP:}\quad & B_i\,\alpha \;\land\; \neg B_i\,(\mathit{Bif}_j\,\alpha \lor \mathit{Uif}_j\,\alpha) \\
\text{RE:}\quad & B_j\,\alpha \\
\text{where}\quad & \alpha = I_i\,\mathit{Done}(\langle j, act\rangle, \phi)
\end{align*}
\end{nbdef}

How to read it, symbol by symbol (the logic is the BDI modal logic of FIPA-SL):
\begin{itemize}[nosep]
  \item $\langle i, X\rangle$: agent $i$ performs communicative act $X$;
  \item $B_i\,p$: $i$ \emph{believes} $p$; $I_i\,p$: $i$ \emph{intends} $p$;
  \item $\mathit{Done}(\langle j, act\rangle, \phi)$: $j$ has performed $act$, and $\phi$ was
        true when it did;
  \item $\mathit{Bif}_j\,\alpha$: $j$ knows whether $\alpha$ (believes $\alpha$ or believes
        $\neg\alpha$); $\mathit{Uif}_j\,\alpha$: $j$ is uncertain about $\alpha$;
  \item \textbf{FP} (\emph{feasibility precondition}): what must hold \emph{for the sender}
        to perform the act sincerely: $i$ believes $\alpha$, and $i$ does not believe that
        $j$ already knows or is uncertain about $\alpha$ (no point in telling otherwise);
  \item \textbf{RE} (\emph{rational effect}): what the sender \emph{wants to achieve}: that
        $j$ comes to believe $\alpha$. The receiver is not forced to adopt it: it is the
        expected effect, not a guaranteed one.
\end{itemize}
In words: \emph{accepting $j$'s proposal means informing $j$ that $i$ intends that $j$ performs
the action under condition $\phi$}. Every FIPA act is defined this way, in terms of mental
states: this is the strong, BDI notion of agency applied to communication.

\begin{nbexample}{accept-proposal (slide 18)}
Agent $i$ informs $j$ that it accepts an offer from $j$ to stream a given multimedia title to
channel 19 when the customer is ready. Agent $i$ will inform $j$ of this fact when
appropriate.
\begin{minted}[linenos=false]{text}
(accept-proposal
  :sender (agent-identifier :name i)
  :receiver (set (agent-identifier :name j))
  :in-reply-to bid089
  :content
    "((action (agent-identifier :name j)
              (stream-content movie1234 19))
      (B (agent-identifier :name j)
         (ready customer78)))"
  :language fipa-sl)
\end{minted}
The content pairs the action (\code{stream-content movie1234 19}, by \code{j}) with the
condition (\code{j} believes that \code{customer78} is ready). \code{:in-reply-to bid089}
links the acceptance to \code{j}'s earlier \code{propose} message.
\end{nbexample}
